\scanpage{26}
\begin{center}$-\ 33\ -$\end{center}
\hfill n$^\circ$ 371

désormais supposées associatives et ayant un \uline{élément unité}.

PROPOSITION 5.--- \uline{Soit} $(E_i)_{i\in I}$ \uline{une famille finie d'algèbres sur un anneau commutatif} $A$, \uline{et pour chaque} $i\in I$, \uline{soit} $e_i$ \uline{l'élément unité de} $E_i$. \uline{Pour chaque} $i\in I$, \uline{soit} $u_i:E_i\to E=\bigotimes_{i\in I}E_i$ \uline{l'application} $A$-\uline{linéaire qui à tout} $x_i\in E_i$ \uline{fait correspondre l'élément} $u_i(x_i)=\bigotimes_{j\in I}x'_j$ \uline{avec} $x'_i=x_i$ \uline{et} $x'_j=e_j$ \uline{pour} $j\ne i$.

(i) \uline{Les} $u_i$ \uline{sont des homomorphismes de} $A$-\uline{algèbres tels que, pour} $i\ne j$, $u_i(x_i)$ \uline{et} $u_j(x_j)$ \uline{soient permutables dans} $E$ \uline{quels que soient} $x_i\in E_i$ \uline{et} $x_j\in E_j$, \uline{et la} $A$-\uline{algèbre} $E$ \uline{est engendrée par la réunion des sous-algèbres} $u_i(E_i)$.

(ii) \uline{Soient} $F$ \uline{une} $A$-\uline{algèbre, et pour tout} $i\in I$, \uline{soit} $v_i:E_i\to F$ \uline{un homomorphisme de} $A$-\uline{algèbres, les} $v_i$ \uline{étant tels que pour} $i\ne j$, $v_i(x_i)$ \uline{et} $v_j(x_j)$ \uline{soient permutables dans} $F$ \uline{quels que soient} $x_i\in E_i$ \uline{et} $x_j\in E_j$. \uline{Alors il existe un homomorphisme de} $A$-\uline{algèbres} $w:E\to F$ \uline{et un seul tel que l'on ait} $v_i=w\circ u_i$ \uline{pour tout} $i\in I$.

(i) Le fait que $u_i$ est un homomorphisme d'algèbres résulte aussitôt de la définition (1) de la multiplication dans $E$. En outre, si $i\ne j$, $x_i\in E_i$, $x_j\in E_j$, on a $u_i(x_i)=\bigotimes_k x'_k$, $u_j(x_j)=\bigotimes_k x''_k$ avec $x'_k=e_k$ pour $k\ne i$, $x'_i=x_i$, $x''_k=e_k$ pour $k\ne j$, $x''_j=x_j$ ; il est clair que cela entraîne $x'_kx''_k=x''_kx'_k$ quels que soit $k\in I$, d'où la seconde assertion. Enfin, on peut écrire pour $(x_i)\in\prod_{i\in I}E_i$, $\bigotimes_{i\in I}x_i=\prod_{i\in I}u_i(x_i)$, le second produit ayant un sens puisqu'il s'agit d'éléments permutables, d'où la troisième assertion.

(ii) L'hypothèse entraîne que pour tout $(x_i)\in\prod_i E_i$ le produit $\prod_{i\in I}v_i(x_i)$ a un sens indépendamment de toute structure d'ordre sur $I$.

L'application $(x_i)\to\prod_i v_i(x_i)$ de $\prod_{i\in I}E_i$ dans $F$ est évidemment $A$-multilinéaire, et il existe donc une application $A$-linéaire et une seule $w:E\to F$ telle que $w(\bigotimes_i x_i)=\prod_i v_i(x_i)$ ; $\unsure{v|w}$ est un homomorphisme d'\uline{algèbres}, car en vertu de l'hypothèse, on a
\mspage{25}
